% Katja Seidel: From the early Common Market to the crises of the 1960s – Reading Summary
% Introduction to the European Union, Week 3
% Compile with: pdflatex seidel-common-market-1960s-reading-summary.tex

\documentclass[11pt,a4paper]{article}
\usepackage[utf8]{inputenc}
\usepackage[T1]{fontenc}
\usepackage[margin=2.5cm]{geometry}
\usepackage{booktabs}
\usepackage{enumitem}
\setlist{nosep,leftmargin=*}
\usepackage{parskip}
\usepackage{microtype}
\usepackage{hyperref}
\usepackage{xcolor}
\definecolor{eublue}{RGB}{0,51,153}

\hyphenation{Spitzenkandidat Spitzenkandidaten}
\hypersetup{
  colorlinks=true,
  linkcolor=eublue,
  urlcolor=eublue,
  pdftitle={Seidel – Common Market 1960s},
  pdfauthor={Introduction to the EU, UCM}
}

\begin{document}

\title{From the early Common Market to the crises of the 1960s, 1958--1968\\Chapter from \textit{Reinventing Europe}}
\author{Introduction to the European Union\\Universidad Complutense de Madrid\\Week 3}
\date{}
\maketitle
\vspace{-1em}
\hrule
\vspace{1em}

\section*{Rhetorical Framework (Ethos, Logos, Pathos)}

\begin{table}[h]
\centering
\begin{tabular}{@{}lll@{}}
\toprule
\textbf{Appeal} & \textbf{Meaning} & \textbf{Seidel's Use} \\
\midrule
\textbf{Ethos} & Credibility, authority, trust & Institutional history; scholarly synthesis; balanced account of crises \\
\textbf{Logos} & Logic, structure, cause-and-effect & Chronological narrative; institutional development; economic vs.\ political outcomes \\
\textbf{Pathos} & Emotion, stakes, human dimension & Crises, stalemate, compromise; tension between federalism and intergovernmentalism \\
\bottomrule
\end{tabular}
\end{table}

\textbf{Key insight:} Seidel emphasizes \textit{logos} (institutions, timelines, outcomes) and \textit{ethos} (authoritative institutional history). \textit{Pathos} appears in the drama of crises---Empty Chair, de Gaulle's resistance, Luxembourg Compromise---and in the contrast between economic success and political deadlock.

\section*{Bibliographic Information}

\begin{itemize}
\item \textbf{Author:} Katja Seidel
\item \textbf{Title:} From the early Common Market to the crises of the 1960s, 1958--1968
\item \textbf{Source:} In LEUCHT, Brigitte; SEIDEL, Katja \& WARLOUZET, Laurent, \textit{Reinventing Europe. The History of the European Union, 1945 to the present}
\item \textbf{Publisher:} Bloomsbury, London
\item \textbf{Year:} 2023
\end{itemize}

\section*{Summary \& Key Points}

\begin{itemize}
\item The period 1958--1968 saw the early development of the Common Market and the establishment of key EEC institutions.
\item Customs union was achieved ahead of schedule in 1968, marking a major success for the European Community.
\item The 1960s were marked by political crises, particularly the Empty Chair Crisis (1965--66) and the Luxembourg Compromise.
\item De Gaulle's leadership shaped this period, with his resignation in 1969 marking a turning point.
\item The CAP was established as a key policy, though controversial.
\item This period established the institutional framework that would guide European integration for decades.
\end{itemize}

\section*{Main Arguments}

\subsection*{I. Institutional Development}

The early years saw the establishment and consolidation of the four main EEC institutions: the Commission, Council of Ministers, Parliament, and Court of Justice. Each institution developed its role and powers during this period.

\subsection*{II. Economic Integration Success}

Despite political challenges, the economic integration project was largely successful. The customs union was achieved ahead of schedule, trade doubled, and living standards improved faster than in the US and UK.

\subsection*{III. Political Crises and Compromises}

The period was marked by significant political crises, particularly the Empty Chair Crisis when France withdrew from Council meetings for six months. The Luxembourg Compromise that resolved it slowed decision-making but kept France in the Community.

\subsection*{IV. De Gaulle's Impact}

Charles de Gaulle's leadership fundamentally shaped this period. His vision of intergovernmental cooperation clashed with more federalist approaches, leading to tensions but also ensuring French participation in the Community.

\section*{Context \& Analysis}

\textbf{Strengths:} Provides detailed account of institutional development and economic achievements during a crucial period.

\textbf{Weaknesses:} May focus more on elite actors and institutions than on social or popular perspectives. Connects to Krawatzek \& Pestel's critique---top-down bias.

\textbf{Questions Raised:} How did ordinary Europeans experience this period? What alternative paths were possible?

\section*{Relevance to Course}

\begin{itemize}
\item \textbf{Ethos:} Essential for understanding early EEC development; authoritative institutional history.
\item \textbf{Logos:} Empty Chair Crisis, Luxembourg Compromise---cause-and-effect, institutional logic.
\item \textbf{Pathos:} Crises and compromises---stakes of integration, national vs.\ supranational.
\item \textbf{All three:} Context for first enlargement (1973); economic success despite political stalemate.
\end{itemize}

\section*{Discussion Questions}

\begin{itemize}
\item Why was the project stalled and how did it get going again? (Ethos: who had power? Logos: what broke the deadlock? Pathos: stakes of stalemate)
\item Significance of Luxembourg Compromise? (Ethos: who won, who lost? Logos: institutional effects; Pathos: compromise as relief or frustration)
\item How did de Gaulle's vision differ? (Ethos: national vs.\ federal voices; Logos: structural tension; Pathos: clash of ideals)
\item Achievements and limitations of first decade? (Ethos: whose success? Logos: evidence, counterevidence; Pathos: hope vs.\ constraint)
\end{itemize}

\vspace{1em}
\noindent\footnotesize\textit{Reference:} Seidel, Katja. ``From the early Common Market to the crises of the 1960s, 1958--1968.'' In LEUCHT, Brigitte; SEIDEL, Katja \& WARLOUZET, Laurent, \textit{Reinventing Europe}, London: Bloomsbury, 2023.

\end{document}
